% =====================================================================
% CHAPTER 1: INTRODUCTION
% =====================================================================
\chapter{Introduction}
\label{ch:introduction}

\section{Problem Statement}

Financial fraud detection is a critical challenge facing the global banking and financial services industry. With the rise of digital banking and online transactions, fraudsters have developed increasingly sophisticated methods to circumvent traditional detection systems. According to industry estimates, financial fraud costs the global economy hundreds of billions of dollars annually \cite{singh2024comprehensive}, with credit card fraud alone accounting for over \$30 billion in losses each year. The fundamental challenge lies in identifying the small fraction of fraudulent transactions (typically less than 1\% of all transactions) among millions of legitimate ones, a problem characterized by extreme class imbalance, evolving fraud patterns, and the need for real-time detection with minimal false alarms.

Traditional fraud detection systems rely on rule-based approaches and tabular machine learning models---which, despite strong performance on structured data \cite{grinsztajn2022tree}---that treat each transaction independently. These methods fail to capture the relational patterns that are often the strongest indicators of organized fraud \cite{pourhabibi2022graph}: shared devices between multiple applicants, coordinated application bursts from the same geographic area, and overlapping behavioral profiles that emerge only when the connections between entities are examined. Moreover, even the best classical machine learning models exhibit misleadingly high accuracy metrics that mask devastatingly poor fraud detection: Ahmed and Shamsuddin \cite{ahmed2022comparative} reported decision tree accuracy of 97.17\% on banking fraud data, yet Dong \cite{dong2025enhanced} demonstrated that Random Forest achieves only 0.16\% fraud recall despite 98.93\% overall accuracy, confirming that conventional metrics are inadequate under extreme class imbalance. Single-paradigm deep learning approaches---whether based solely on graph neural networks, standard feedforward architectures, or isolated anomaly detection models---suffer from complementary limitations. Purely graph-based methods depend heavily on graph construction quality and struggle with temporal dynamics \cite{vanderhulst2024graph}; standard neural networks lack relational reasoning; and unsupervised anomaly detectors cannot exploit labeled fraud patterns effectively. Standalone autoencoders achieve near-random performance using reconstruction error alone, with Dong \cite{dong2025enhanced} reporting a DAE AUC of only 0.5210 without supervised stacking. Spiking neural networks face calibration instability and temporal generalization failure on distributional shift scenarios \cite{ahmed2025bioinspired}. No single learning paradigm fully addresses the multifaceted nature of financial fraud, motivating the exploration of a multi-paradigm approach that combines the strengths of graph-based, neuromorphic, and unsupervised deep learning.

\section{Motivation and Significance}

The motivation for this research stems from five key observations. First, financial fraud is inherently relational: fraud rings operate through shared resources (devices, IP addresses, email domains) that create detectable network patterns when represented as a graph. Second, graph neural networks (GNNs) have demonstrated state-of-the-art performance in domains where relational structure is informative \cite{pourhabibi2022graph,cheng2024graph}, including social network analysis, drug discovery, and recommendation systems. Third, the application of GNNs to financial fraud detection is still an emerging field, with significant gaps in understanding the relative strengths of different GNN architectures, the impact of graph construction strategies, and the robustness of these models under distributional shift. Fourth, spiking neural networks (SNNs) offer a biologically inspired computing paradigm with inherent advantages for temporal pattern recognition and event-driven processing \cite{nasif2026reinforcement,ahmed2025bioinspired}, yet their application to financial fraud detection remains virtually unexplored and faces calibration challenges under extreme class imbalance. Fifth, unsupervised autoencoders---particularly Variational Autoencoders (VAEs) \cite{kingma2014auto} and Denoising Autoencoders (DAEs) \cite{vincent2010stacked}---provide a powerful anomaly detection framework that does not require labeled fraud examples, making them highly suitable for domains where fraud labels are scarce or delayed. However, standalone autoencoders achieve near-random anomaly detection performance on highly imbalanced tabular data \cite{dong2025enhanced,obushyni2025vae}, necessitating hybrid stacking approaches to achieve competitive results.

The significance of this work lies in its potential to advance both the theoretical understanding and practical deployment of multi-paradigm fraud detection systems. By systematically comparing and combining three fundamentally different deep learning paradigms---graph-based relational reasoning (GNN), neuromorphic temporal processing (SNN), and unsupervised anomaly detection (Autoencoder)---this thesis provides actionable insights for practitioners seeking to deploy robust, multi-faceted fraud detection in production environments. Each paradigm addresses a distinct aspect of the fraud detection challenge: GNNs exploit relational structure, SNNs capture temporal dynamics through spike-based processing, and autoencoders enable label-independent anomaly scoring.

\section{Research Gaps, Questions, and Objectives}

\subsection{Research Gaps}

Despite the growing interest in deep learning-based fraud detection, several critical gaps remain in the literature:

\begin{enumerate}[label=\arabic*.]
\item \textbf{Graph construction methodology:} Existing works often treat graph construction as a preprocessing afterthought, without systematic analysis of how node type selection, feature bucketing strategies, and edge construction policies affect downstream model performance \cite{islam2023effects}.
\item \textbf{GNN architecture comparison:} Most studies evaluate a single GNN architecture without systematic comparison across fundamentally different paradigms (spatial vs.\ quantum-inspired vs.\ spectral approaches).
\item \textbf{Hyperparameter sensitivity:} The sensitivity of spectral GNN architectures to hyperparameters such as wavelet order and graph splitting configurations is poorly understood.
\item \textbf{Fairness and robustness:} The performance stability of GNN-based fraud detectors under distributional shift and across different data versions has received limited attention \cite{chouldechova2017fair,kozodoi2022fairness}.
\item \textbf{Class imbalance handling:} The effectiveness of resampling strategies \cite{chawla2002smote,he2008adasyn} in the context of graph-based learning, where neighborhood structure is critical, remains unexplored.
\item \textbf{SNNs for fraud detection:} The application of spiking neural networks to financial fraud detection is virtually unexplored \cite{nasif2026reinforcement}. Key questions remain regarding the suitability of LIF neuron models, rate coding strategies, surrogate gradient descent training, and ensemble methods for the extreme class imbalance and high-dimensional feature spaces characteristic of fraud data.
\item \textbf{Autoencoder ensembles for fraud:} While individual autoencoder architectures (VAE \cite{kingma2014auto}, DAE \cite{vincent2010stacked}) have been applied to anomaly detection, standalone autoencoders achieve near-random performance on imbalanced tabular fraud data, with DAEs reporting AUC as low as 0.5210 \cite{dong2025enhanced} and standard autoencoders achieving AUC of approximately 0.55 \cite{obushyni2025vae}. The systematic combination of complementary autoencoder paradigms through score-level stacking with gradient-boosted meta-learners \cite{chen2016xgboost,ke2017lightgbm} has not been fully investigated for financial fraud, and the resampling paradox---where synthetic oversampling degrades autoencoder performance---requires further investigation \cite{dong2025enhanced}. The optimal design of swap noise corruption, latent space dimensionality, and evidence lower bound optimization \cite{farahat2025objective} for fraud-specific anomaly detection remains open.
\item \textbf{Cross-paradigm comparison:} No existing work systematically compares GNN, SNN, and autoencoder-based approaches on the same fraud detection task \cite{pourhabibi2022graph,cheng2024graph}, limiting practitioners' ability to select or combine paradigms based on their complementary strengths and weaknesses.
\end{enumerate}

\subsection{Research Questions}

\begin{enumerate}[label=RQ\arabic*.]
\item How can tabular financial fraud data be effectively transformed into a heterogeneous graph representation that captures relational fraud patterns?
\item Which GNN architecture paradigm (spatial aggregation, quantum-inspired, or spectral splitting) provides the best trade-off between fraud detection performance and computational efficiency?
\item How sensitive is the Split GNN architecture to its key hyperparameters (wavelet order $C$, split point $K$, edge loss weight $\gamma_{\text{edge}}$), and what is the optimal configuration?
\item How robust are GNN-based fraud detectors to distributional shifts across different data versions, and what are the failure modes?
\item Does resampling-based class balancing improve graph-based fraud detection compared to loss-based class weighting \cite{farahat2025objective}?
\item Can spiking neural networks with Leaky Integrate-and-Fire neurons and rate coding achieve competitive fraud detection performance \cite{nasif2026reinforcement}, and how does surrogate gradient descent training perform under extreme class imbalance?
\item How do Variational Autoencoders \cite{kingma2014auto} and Denoising Autoencoders \cite{vincent2010stacked} compare as unsupervised anomaly detectors for fraud, and does score-level stacking with XGBoost/LightGBM \cite{chen2016xgboost,ke2017lightgbm} improve detection over individual models?
\item What are the complementary strengths and weaknesses of the GNN, SNN, and Autoencoder pipelines, and how can they be effectively combined or selected for production deployment?
\end{enumerate}

\subsection{Research Objectives}

\begin{enumerate}[label=O\arabic*.]
\item Design and implement a graph construction pipeline (\textsc{FraudGraphBuilder}) that transforms tabular fraud data into a heterogeneous graph with 17 node types.
\item Implement and evaluate three GNN architectures: GraphSAGE, Quantum GNN, and Split GNN.
\item Conduct a comprehensive 45-experiment hyperparameter tuning study for the Split GNN.
\item Evaluate model fairness and robustness across six data distribution variants.
\item Compare class weighting and resampling strategies for handling extreme class imbalance in graph-based fraud detection.
\item Design and implement an SNN pipeline using Leaky Integrate-and-Fire neurons, rate coding, and surrogate gradient descent, and evaluate ensemble methods for improving SNN-based fraud detection.
\item Design and implement an Autoencoder pipeline combining a Variational Autoencoder (VAE) and a Denoising Autoencoder (DAE) with score-level stacking using XGBoost and LightGBM meta-learners.
\item Conduct a systematic cross-pipeline comparison of GNN, SNN, and Autoencoder approaches, identifying complementary strengths and providing practical deployment recommendations.
\end{enumerate}

\section{Proposed Solution Overview}

This thesis proposes a comprehensive multi-paradigm fraud detection framework that addresses the identified research gaps. The framework consists of five interconnected components:

\textbf{Graph Construction:} The \textsc{FraudGraphBuilder} transforms flat tabular applicant data into a heterogeneous graph with 17 distinct node types (categorical, bucketed continuous, and binary/flag), creating shared attribute nodes that enable the detection of relational fraud patterns through indirect user connections.

\textbf{GNN Pipeline:} Three GNN architectures are implemented and compared \cite{cheng2024graph}. GraphSAGE provides a strong classical baseline through neighborhood sampling and aggregation. The Quantum GNN explores quantum-inspired feature transformations through angle encoding and variational quantum circuits. The Split GNN introduces a novel graph-partitioning mechanism that combines edge classification with dual-branch Beta wavelet filtering for frequency-domain feature extraction. A 45-experiment grid search over wavelet order ($C \in \{1,2,3\}$), split point ($K \in \{-1,0,1,2\}$), and edge loss weight ($\gamma_{\text{edge}} \in \{0.2, 0.4, 0.6, 0.8, 1.0\}$) reveals the sensitivity and optimal configuration of the Split GNN architecture.

\textbf{SNN Pipeline:} A spiking neural network pipeline employs Leaky Integrate-and-Fire (LIF) neurons with rate coding to process tabular fraud features as temporal spike trains \cite{nasif2026reinforcement}. Surrogate gradient descent enables effective backpropagation through the non-differentiable spike threshold, while ensemble methods combining multiple SNN configurations enhance detection stability. The SNN paradigm offers inherent advantages for event-driven transaction processing on neuromorphic hardware, though software-based simulation introduces computational overhead.

\textbf{Autoencoder Pipeline:} An unsupervised anomaly detection pipeline combines a Variational Autoencoder (VAE) \cite{kingma2014auto}, which models the latent distribution of legitimate transactions using the ELBO objective with the reparameterization trick, and a Denoising Autoencoder (DAE) \cite{vincent2010stacked}, which learns robust representations by reconstructing inputs corrupted with swap noise. Score-level stacking of VAE and DAE anomaly scores with gradient-boosted classifiers (XGBoost \cite{chen2016xgboost} and LightGBM \cite{ke2017lightgbm}) produces a powerful ensemble that effectively separates fraudulent from legitimate patterns without requiring labeled training data.

\textbf{Cross-Pipeline Comparison and Robustness Evaluation:} Cross-version validation across six data distributions quantifies model stability under distributional shift, a controlled resampling experiment evaluates the effectiveness of ADASYN-based class balancing \cite{he2008adasyn} versus inverse-frequency weighting, and a systematic cross-pipeline comparison identifies the complementary strengths and weaknesses of each paradigm for practical deployment.

\section{Thesis Organisation}

The remainder of this thesis is organised as follows:

\textbf{Chapter~\ref{ch:litreview}} presents a comprehensive literature review covering the domain background of financial fraud detection, the state-of-the-art in software solutions, a detailed survey of traditional, classical machine learning, graph neural network, spiking neural network, and autoencoder approaches, a comparative analysis across paradigms, and the identification of research gaps that motivate this work.

\textbf{Chapter~\ref{ch:requirements}} presents the analysis and requirements engineering, including stakeholder analysis, functional and non-functional requirements, and data requirements for the proposed system.

\textbf{Chapter~\ref{ch:architecture}} describes the system architecture and design, including the high-level system architecture, component design, the AI pipeline architecture, and the technology stack selection.

\textbf{Chapter~\ref{ch:methodology}} details the ML methodology and implementation, covering data preprocessing techniques, model selection and justification, feature engineering and selection, hyperparameter optimisation strategy, and software implementation details.

\textbf{Chapter~\ref{ch:gnn}} presents the GNN pipeline, including the design and evaluation of GraphSAGE, Quantum GNN, and Split GNN architectures, hyperparameter sensitivity analysis, and cross-version fairness evaluation.

\textbf{Chapter~\ref{ch:snn}} presents the SNN pipeline, covering the design and implementation of Leaky Integrate-and-Fire neurons, rate coding strategies, surrogate gradient descent training, and ensemble methods for spiking neural network-based fraud detection.

\textbf{Chapter~\ref{ch:autoencoder}} presents the Autoencoder pipeline, including the design and evaluation of Variational Autoencoder and Denoising Autoencoder models, anomaly score computation, and score-level stacking with XGBoost/LightGBM meta-learners.

\textbf{Chapter~\ref{ch:comparison}} provides a cross-pipeline comparison, analysing the complementary strengths and weaknesses of the GNN, SNN, and Autoencoder approaches, and offering practical recommendations for multi-paradigm deployment.

\textbf{Chapter~\ref{ch:conclusion}} concludes the thesis with a summary of contributions, critical reflection, and recommendations for future research.

\textbf{Appendices} provide the AI Ethics \& Bias Statement, Reproducibility Report, Generative AI Disclosure, Extended Hyperparameter Configurations, Denoising Autoencoder Architecture and Objective Formulations, Spiking Neural Network Mathematical Derivations, and Graph Neural Network Algorithmic Procedures.
